\documentclass{article}
\usepackage[preprint]{neurips_2025}
\usepackage[T1]{fontenc}
\usepackage{amsmath,amssymb,booktabs,tabularx}
\usepackage{microtype}
\usepackage[hidelinks]{hyperref}
\usepackage{url}
\title{Asymmetric Process-Reward Adaptation\\for Small-Model Mathematical RL}
\author{
Hongyi Ding \\
\And
Jasmine Jiang \\
\And
Yang Zeng \\
}
\newcommand{\E}{\mathbb{E}}
\newcommand{\sg}{\operatorname{sg}}
\newcolumntype{L}{>{\raggedright\arraybackslash}X}
\begin{document}
\maketitle
\begin{abstract}
Final-answer correctness can provide noisy supervision for process reward models (PRMs). We test one loss change: omit positive outcome supervision during offline PRM adaptation while retaining public step labels. Four matched GRPO experiments measure policy-only mathematical accuracy and reward quality under a small-model budget.
\end{abstract}

\section{Research problem and importance}
Some answer-correct solutions contain reasoning errors \citep{zheng2025}. Their positive outcome labels can teach a PRM to reward faulty reasoning. We ask: \textbf{does omitting this supervision improve the policy's mathematical accuracy?} We hypothesize that reducing false certification can outweigh losing useful positive supervision; downstream experiments must establish whether this trade-off is favorable.

\section{Related work, limitations, and contribution}
PRM800K supplies step labels \citep{lightman2023}; PRIME updates implicit rewards online \citep{cui2025}; PURE addresses reward hacking through min-form credit assignment \citep{cheng2025}. LCA learns PRMs from outcomes with softmax-weighted pooling and already studies one-sided label noise \citep{jia2026}. We use its pooling but test a narrower question: \textbf{whether omitting positive outcome supervision, while retaining real step labels, improves downstream GRPO}. The symmetric control is LCA-inspired, not a full reproduction. Our contribution is this controlled supervision change and its policy-level evaluation.

\section{Proposed method}
\textbf{Adapt once.} Generate four solutions per problem on 512 training problems using the initial policy. Retain complete, unambiguously verified solutions. For step-boundary correctness scores $p_{\phi,t}$, set $e_{\phi,t}=1-p_{\phi,t}$. Following LCA,
\begin{equation}
 w_t=\frac{\exp(e_{\phi,t}/T)}{\sum_j\exp(e_{\phi,j}/T)},\qquad
 E_\phi(\tau)=\sum_t\sg(w_t)e_{\phi,t},\qquad T=0.2 .
\end{equation}
Here $\sg$ stops gradients through weights. Let $L_{\rm fail}=\E_{\tau^-}[-\log E_\phi(\tau^-)]$ and $L_{\rm success}=\E_{\tau^+}[-\log(1-E_\phi(\tau^+))]$, using separate means over failures/successes. A shared set of 2,048 public positive/negative step labels supplies mean binary cross-entropy $L_{\rm step}$:
\begin{equation}
 L_c=L_{\rm fail}+cL_{\rm success}+L_{\rm step},\qquad
 c_C=1,\quad \boxed{c_D=0\ \text{(ours)}} .
\end{equation}
Failures provide weak evidence of an error, without its location. Real step labels discourage blanket rejection; they do not guarantee accurate scores or downstream gains.

\newpage
\noindent\textbf{Policy training (continued).} For B--D, use the frozen PRM to compute
\begin{equation}
 R_i=y_i+0.2\min_t p_{\phi,t}(\tau_i),\qquad
 \hat A_i=\frac{R_i-\bar R}{\sigma_R+10^{-4}} .
\end{equation}
Here $y_i\in\{0,1\}$ is answer correctness; $\bar R,\sigma_R$ are the group mean/std. A uses $R_i=y_i$. The detached sequence advantage applies to all completion tokens in clipped GRPO \citep{shao2024}. Correct policy rollouts remain in training.

\section{Four-group experimental design}
\begin{center}
\small
\renewcommand{\arraystretch}{1.12}
\begin{tabularx}{\linewidth}{@{}l L L@{}}
\toprule
\textbf{Group} & \textbf{Reward / PRM preparation} & \textbf{Controlled comparison} \\
\midrule
A & Outcome only; no PRM & Outcome-only GRPO baseline \\
B & Mixed reward; original PRM frozen & B vs.\ A: existing process reward \\
C & Mixed reward; adapt with c = 1, then freeze & C vs.\ B: full adaptation package \\
\textbf{D (ours)} & Mixed reward; adapt with c = 0, then freeze & \textbf{D vs.\ C: proposed supervision rule} \\
\bottomrule
\end{tabularx}
\end{center}
\textbf{Implementation.} Use Qwen2.5-Math-1.5B as policy and H4 Qwen2.5-Math-1.5B-Instruct-PRM-0.2 as PRM \citep{h4prm}; LoRA ranks are 16/8. Preserve native step formatting and freeze PRMs throughout RL. TRL/PEFT uses \texttt{loss\_type="grpo"} and sums rewards before group normalization \citep{trl}. Truncated/unparseable outputs receive zero reward. Pin all configurations before formal runs.

\textbf{Data.} Reserve 500 MATH training problems for validation; use the remainder for RL/buffer generation \citep{hendrycks2021}. Deduplicate new training data/anchors against validation, tests, and ProcessBench. PRM800K includes 4,500 original MATH test problems in training \citep{prmdata}; filtering new data cannot undo inherited PRM exposure. Use \textbf{MATH500 as primary test after checkpoint-corpus overlap audit}, GSM8K for transfer \citep{cobbe2021}, and full MATH only as exposure-aware supplementary evaluation. Unknown pretraining exposure remains a limitation.

\textbf{Controls and budget.} C/D share initialization, batches, optimizer, and updates; only $c$ changes, without re-averaging losses. All arms share policy initialization, prompt order, decoding, and tuning effort. Start with eight rollouts/question, temperature 1.0, and a 2,048-token cap. A one-/two-RTX-5090 pilot fixes a common budget: provisionally 6M generated tokens/run, four groups, three paired seeds (72M tokens). Report actual updates, prompt counts, length, truncation, and total GPU hours including all data generation, adaptation, scoring, training, and evaluation.

\textbf{Policy evaluation.} Evaluate the initial policy and final fixed-budget checkpoints with identical greedy decoding, without PRM reranking or voting. \textbf{D--C MATH500 accuracy is the primary comparison}; report per-seed differences, mean/std, and 95\% paired question-bootstrap intervals. D vs.\ A/B tests practical value; small or inconsistent gains remain inconclusive.

\textbf{Reward diagnostic.} On fixed ProcessBench MATH solutions, report first-error accuracy, correct-chain acceptance, and their harmonic mean at threshold 0.5. Within answer-correct solutions, valid-chain AUROC tests discrimination beyond a global score shift. Report overlaps and class counts; no diagnostic is used for tuning. Reward hacking and policy distribution shift remain risks.

\section{Three-person plan and tentative timeline}
Hongyi Ding leads GRPO and GPU profiling; Jasmine Jiang leads PRM data and adaptation; Yang Zeng leads split audits, evaluation, and statistics. All three write the report.
\begin{center}
\small
\renewcommand{\arraystretch}{1.1}
\begin{tabularx}{\linewidth}{@{}l L@{}}
\toprule
\textbf{Period} & \textbf{Deliverable and responsibility} \\
\midrule
Week 1 & Hongyi: 48-hour resource/loss pilot; Jasmine: buffer and PRMs; Yang: split audit. Proposal due Oct.\ 21, 11:59\,pm PT. \\
Week 2 & Hongyi/Jasmine: first-seed A--D runs; Yang: verify rewards, accuracy, and cost logs. \\
Week 3 & Hongyi/Jasmine: remaining seeds; Yang: paired comparisons and error analysis. \\
Week 4 & All: final evaluation, reproducible code, figures, and report. \\
\bottomrule
\end{tabularx}
\end{center}
\clearpage
\begin{thebibliography}{11}
\bibitem[Shao et al.(2024)]{shao2024}
Zhihong Shao et al. DeepSeekMath: Pushing the Limits of Mathematical Reasoning in Open Language Models. 2024.
\url{https://arxiv.org/abs/2402.03300v3}.

\bibitem[Zheng et al.(2025)]{zheng2025}
Chujie Zheng et al. ProcessBench: Identifying Process Errors in Mathematical Reasoning. ACL, 2025.
\url{https://arxiv.org/abs/2412.06559v4}.

\bibitem[Lightman et al.(2023)]{lightman2023}
Hunter Lightman et al. Let's Verify Step by Step. 2023.
\url{https://arxiv.org/abs/2305.20050}.

\bibitem[Cui et al.(2025)]{cui2025}
Ganqu Cui et al. Process Reinforcement through Implicit Rewards. 2025.
\url{https://arxiv.org/abs/2502.01456v2}.

\bibitem[Cheng et al.(2025)]{cheng2025}
Jie Cheng et al. Stop Summation: Min-Form Credit Assignment Is All Process Reward Model Needs for Reasoning. NeurIPS, 2025.
\url{https://arxiv.org/abs/2504.15275v3}.

\bibitem[Jia et al.(2026)]{jia2026}
Tianyu Jia et al. The Weakest Link Tells It All: Outcome-Supervised Process Reward Modeling via Learnable Credit Assignment. EMNLP, 2026. Camera-ready version.
\url{https://arxiv.org/abs/2606.27739v2}.

\bibitem[Hendrycks et al.(2021)]{hendrycks2021}
Dan Hendrycks et al. Measuring Mathematical Problem Solving With the MATH Dataset. NeurIPS, 2021.
\url{https://arxiv.org/abs/2103.03874v2}.

\bibitem[Cobbe et al.(2021)]{cobbe2021}
Karl Cobbe et al. Training Verifiers to Solve Math Word Problems. 2021.
\url{https://arxiv.org/abs/2110.14168v2}.

\bibitem[Hugging Face H4(2025)]{h4prm}
Hugging Face H4. Qwen2.5-Math-1.5B-Instruct-PRM-0.2. 2025.
\href{https://huggingface.co/HuggingFaceH4/Qwen2.5-Math-1.5B-Instruct-PRM-0.2}{Official model card and scoring interface}. Accessed Oct.\ 19, 2026.

\bibitem[Hugging Face(2026)]{trl}
Hugging Face. TRL: \href{https://huggingface.co/docs/trl/grpo_trainer}{GRPO Trainer documentation}. Accessed Oct.\ 19, 2026.

\bibitem[OpenAI(2023)]{prmdata}
OpenAI. PRM800K dataset repository: MATH splits and answer grading. 2023. Accessed Oct.\ 19, 2026.
\url{https://github.com/openai/prm800k}.
\end{thebibliography}
\end{document}
